\documentclass[11pt,a4paper]{article}
\usepackage[T1]{fontenc}
\usepackage[utf8]{inputenc}
\usepackage{amsmath,amsthm}
\usepackage{newpxtext,newpxmath}
\usepackage[margin=25mm,headheight=15pt]{geometry}
\usepackage{microtype,mathtools,booktabs,array,longtable,xcolor}
\usepackage{enumitem,fancyhdr,xurl,hyperref,cleveref}
\usepackage{aliascnt}
\hypersetup{colorlinks=true,linkcolor=blue!40!black,citecolor=blue!40!black,
 urlcolor=blue!40!black,pdftitle={Quadratic Renyi Growth in Geometric Convolution Channels},
 pdfauthor={Research article prepared for Vladimir Reshetnikov}}
\pagestyle{fancy}\fancyhf{}
\fancyhead[L]{\small Quadratic R\'enyi growth}
\fancyhead[R]{\small Geometric convolution channels}
\fancyfoot[L]{\small 30 September 2026}\fancyfoot[R]{\thepage}
\setlength{\parskip}{4pt}\setlength{\parindent}{1.1em}
\setlist{nosep,leftmargin=*}
\numberwithin{equation}{section}
\newtheorem{theorem}{Theorem}[section]
\newaliascnt{lemma}{theorem}
\newtheorem{lemma}[lemma]{Lemma}
\aliascntresetthe{lemma}
\crefname{lemma}{lemma}{lemmas}
\Crefname{lemma}{Lemma}{Lemmas}
\newaliascnt{proposition}{theorem}
\newtheorem{proposition}[proposition]{Proposition}
\aliascntresetthe{proposition}
\crefname{proposition}{proposition}{propositions}
\Crefname{proposition}{Proposition}{Propositions}
\newaliascnt{corollary}{theorem}
\newtheorem{corollary}[corollary]{Corollary}
\aliascntresetthe{corollary}
% ed. (2026-09-30): plural names corrected (delivered: "corollarys", "Corollarys").
\crefname{corollary}{corollary}{corollaries}
\Crefname{corollary}{Corollary}{Corollaries}
\theoremstyle{definition}\newaliascnt{definition}{theorem}
\newtheorem{definition}[definition]{Definition}
\aliascntresetthe{definition}
\crefname{definition}{definition}{definitions}
\Crefname{definition}{Definition}{Definitions}
\newaliascnt{problem}{theorem}
\newtheorem{problem}[problem]{Research question}
\aliascntresetthe{problem}
\crefname{problem}{research question}{research questions}
\Crefname{problem}{Research question}{Research questions}
\theoremstyle{remark}\newaliascnt{remark}{theorem}
\newtheorem{remark}[remark]{Remark}
\aliascntresetthe{remark}
\crefname{remark}{remark}{remarks}
\Crefname{remark}{Remark}{Remarks}
% ed. (2026-09-30): unnumbered environment for ProveIt editorial notes (the theorem counter is unchanged).
\newtheorem*{ednote}{Editorial note (ProveIt, 2026-09-30)}
\newcommand{\E}{\mathbb E}\newcommand{\Prb}{\mathbb P}
\newcommand{\R}{\mathbb R}\newcommand{\N}{\mathbb N}
\newcommand{\dd}{\,\mathrm d}
\newcommand{\KL}{\operatorname{D}_{\mathrm{KL}}}
\newcommand{\Bf}{\mathrm B}
\newcommand{\supp}{\operatorname{supp}}
\newcommand{\Var}{\operatorname{Var}}
\newcommand{\logp}{\log^+}
\newcommand{\eps}{\varepsilon}
\newcommand{\asympc}{\asymp}
\newcommand{\repo}{\texttt{ProveIt}}
\makeatletter\let\inputtable\@@input\makeatother
\title{\LARGE\bfseries Quadratic R\'enyi Growth\\[5pt]
\Large A Universal Linear Correction and Endpoint Selection\\
in Geometric Convolution Channels}
\author{Research article prepared for Vladimir Reshetnikov}
\date{30 September 2026}
\begin{document}
\maketitle
\begin{abstract}
The \repo\ research draft \emph{Beyond Polynomial Diagonals} proves that the
R\'enyi information between a geometric-uniform random series and its first
$m$ terms has a critical order at two, but leaves the supercritical growth
rate open. We resolve that fixed-order growth question. If $0<q<1$ is fixed,
$L=\log(1/q)$, and $\alpha>2$, then
\[
 I_\alpha(S_m;X)=\frac{L}{2}(\alpha-2)m^2+B(\alpha)m
                  +O_{q,\alpha}\bigl((\log(m+1))^2\bigr),
\]
where
\[
 B(\alpha)=\frac{\alpha(\alpha-2)}{\alpha-1}\log\alpha
                  -(\alpha-1)\log(\alpha-1).
\]
In particular, the linear coefficient is independent of $q$ and there is no
$m\log m$ term. We prove the broader result for independent
$\operatorname{Beta}(\theta_0,\theta_1)$ innovations with
$\theta_0,\theta_1\ge1$: both displayed coefficients are multiplied by
$\max\{\theta_0,\theta_1\}$. The endpoint with the larger exponent supplies the
dominant contribution.

The proof combines a positive-simplex estimate for the endpoint density,
a uniform finite-prefix estimate, and a two-variable saddle analysis whose
allocation integral is asymptotically a beta integral. Under the probability
measure tilted by the R\'enyi integrand, a dominant endpoint is approached on
the scale $q^{(\alpha-1)m+o(m)}$, the residual tail on the scale
$q^{(\alpha-2)m+o(m)}$, and the prefix accounts for the fraction
$(\alpha-1)^{-2}$ of the endpoint deficit. We establish large-deviation
statements for these assertions and supply exact algebraic and dyadic-density
checks. These are mathematical proofs with computational diagnostics, not Lean
certificates; literature-wide priority and a uniform critical crossover are
not claimed.
\end{abstract}
\newpage
\tableofcontents
\newpage

\section{The research question and the answer}
\label{sec:introduction}

\subsection{The repository-local open problem}
The information-theoretic research program under
\path{Analysis/FabiusFunction/docs/semi-formalized-research-frontiers/}
studies the statistical content of the finite prefixes of an infinite
geometric convolution. Its \emph{Nuclear Bayesian Operators} draft
\cite{prior} establishes the existence of every finite R\'enyi order at every
finite depth, and proves a transition in the large-depth excess above Shannon
information. In its section ``Further research questions and proposed
routes,'' the problem headed ``Growth above the critical R\'enyi order'' asks
for matching growth estimates when $\alpha>2$, explicitly proposing quadratic
growth in $m$ as a first target. The present article answers that question and
identifies the next coefficient.

This is a specific continuation, not another treatment of the polynomial
singular spectrum, exact null modes, or the prefix-only rate-distortion
obstruction already discussed in that draft. We do not use an unreviewed
repository claim as a black box in the new asymptotic proof. All density and
integral estimates needed below are proved here.

The R\'enyi information throughout is
\begin{equation}
 I_\alpha(S;X)=D_\alpha(P_{S,X}\Vert P_S\otimes P_X).
 \label{eq:information-convention}
\end{equation}
It is not Sibson information, Arimoto information, or a minimized conditional
R\'enyi quantity. These different definitions need not have the same
asymptotics. Our divergence convention is the standard one in
\cite{renyi}.

\subsection{The uniform channel}
Let $U_0,U_1,\ldots$ be independent, uniform on $[0,1]$, and define
\begin{equation}
 X=(1-q)\sum_{j\ge0}q^jU_j,
 \qquad S_m=(1-q)\sum_{j=0}^{m-1}q^jU_j,
 \qquad r=q^m,\quad L=\log(1/q).
 \label{eq:uniform-channel}
\end{equation}
Then $X=S_m+rZ$, where $Z$ is independent of $S_m$ and has the same law as
$X$. The support of $X$ is $[0,1]$ and that of $S_m$ is $[0,1-r]$.
At $q=1/2$ this is the standard geometric-uniform probability construction
behind the Fabius function; its density is the rescaled Rvachev up-function
\cite{arias}. In this article $f$ always denotes the \emph{density}, not the
Fabius cumulative distribution function.

The predecessor uses innovations on $[-1,1]$. Its variables are obtained by
$X_{\mathrm{old}}=2X-1$ and
$S_{m,\mathrm{old}}=2S_m-(1-r)$. Applying these separate invertible affine
maps leaves the divergence in \eqref{eq:information-convention} unchanged,
so the information values and our answer use exactly the same normalization.

\begin{theorem}[The supercritical uniform law]
\label{thm:uniform}
For every fixed $q\in(0,1)$ and $\alpha>2$, as $m\to\infty$,
\begin{equation}
 \boxed{\displaystyle
 I_\alpha(S_m;X)=\frac{\alpha-2}{2}L m^2+B(\alpha)m
                   +O_{q,\alpha}\bigl((\log(m+1))^2\bigr),}
 \label{eq:uniform-main}
\end{equation}
where
\begin{equation}
 B(\alpha)=\frac{\alpha(\alpha-2)}{\alpha-1}\log\alpha
                -(\alpha-1)\log(\alpha-1).
 \label{eq:Balpha}
\end{equation}
Consequently,
\begin{align}
 \lim_{m\to\infty}\frac{I_\alpha(S_m;X)}{m^2}
      &=\frac{\alpha-2}{2}\log(1/q),\label{eq:quadratic-limit}\\
 \lim_{m\to\infty}\frac{I_\alpha(S_m;X)
                  -\frac12(\alpha-2)Lm^2}{m}
      &=B(\alpha).\label{eq:linear-limit}
\end{align}
\end{theorem}
The statement is for \emph{fixed} $q$ and \emph{fixed} $\alpha>2$.
In particular, substituting an $m$-dependent order into
\eqref{eq:uniform-main} is not justified by this theorem.

The Shannon information is exactly $mL$. Thus the open problem phrased in
terms of excess information has the answer
\begin{equation}
 I_\alpha(S_m;X)-mL
   =\frac{\alpha-2}{2}Lm^2+\bigl(B(\alpha)-L\bigr)m
     +O\bigl((\log(m+1))^2\bigr).
 \label{eq:excess-main}
\end{equation}
Notice the subtraction of $L$ from the linear coefficient of the
\emph{excess}. Confusing information with excess information would lose
precisely this term.

\subsection{A universality theorem}
The phenomenon is not restricted to uniform innovations. Let the $U_j$ instead
have the beta density
\begin{equation}
 h(u)=\frac{u^{\theta_0-1}(1-u)^{\theta_1-1}}
                   {\Bf(\theta_0,\theta_1)}\,\boldsymbol1_{(0,1)}(u),
 \qquad \theta_0,\theta_1\ge1,
 \label{eq:beta-innovation}
\end{equation}
and define $X,S_m,r,L$ by the same formulas. Write
$\theta_* =\max\{\theta_0,\theta_1\}$.

\begin{theorem}[Beta-innovation universality and endpoint selection]
\label{thm:beta-main}
For fixed parameters in \eqref{eq:beta-innovation}, $q\in(0,1)$ and
$\alpha>2$,
\begin{equation}
 I_\alpha(S_m;X)
   =\theta_*\left[\frac{\alpha-2}{2}Lm^2+B(\alpha)m\right]
     +O\bigl((\log(m+1))^2\bigr).
 \label{eq:beta-main}
\end{equation}
If $\theta_0>\theta_1$, the lower endpoint dominates the R\'enyi integral;
if $\theta_1>\theta_0$, the upper endpoint dominates. If the exponents are
equal, the channel is symmetric and its two endpoint contributions are equal.
The precise exponential comparison is given in \cref{cor:selection}.
\end{theorem}
The larger exponent means a thinner innovation density near that endpoint.
The high-order likelihood moment is therefore governed by the more strongly
suppressed endpoint, rather than the more probable one. The finite prefix
and the infinite tail share this suppression, and their relative constraints
produce the explicit competition studied below.

\subsection{Scope and status of the contribution}
The geometric series representation, R\'enyi divergence, the beta integral,
Stirling asymptotics, and the Dirichlet simplex distribution are established
tools \cite{arias,renyi,dlmf-beta,dlmf-gamma}. They are not claimed as new.
The contribution is the proved two-term supercritical law, its beta-family
universality, and the associated selection and allocation results for this
channel. The inspected predecessor does not contain these results; its
supercritical question is explicitly left open. A targeted search of the
available primary literature did not establish a prior identical theorem,
but that is not a certification of literature-wide priority.

\section{Exact channel identities and the critical order}
\label{sec:channel}
Put $c=1-q$, and denote the densities of $X$ and $S_m$ by $f$ and $b_m$.
The density of $rZ$ is $r^{-1}f(\cdot/r)$, so
\begin{equation}
 f(x)=\int b_m(s)r^{-1}f((x-s)/r)\dd s.
 \label{eq:convolution}
\end{equation}
The joint density of $(S_m,Z)$ is $b_m(s)f(z)$. For $\alpha>1$, put
$a=\alpha-1$ and define
\begin{align}
 J_{\alpha,m}&=\exp\bigl(a I_\alpha(S_m;X)\bigr),\label{eq:J}\\
 M_{\alpha,m}&=r^aJ_{\alpha,m}
       =\exp\bigl(a(I_\alpha(S_m;X)-mL)\bigr).
 \label{eq:M}
\end{align}
A change of variables $x=s+rz$ gives
\begin{equation}
 M_{\alpha,m}
 =\iint b_m(s)f(z)^\alpha f(s+rz)^{-a}\dd s\dd z.
 \label{eq:master-integral}
\end{equation}
All integrands here and below are nonnegative. Integrals can first be
interpreted in $[0,\infty]$; finiteness will follow from the endpoint estimates.

\begin{lemma}[Conditional form]
\label{lem:conditional}
For $\alpha>1$,
\begin{equation}
 M_{\alpha,m}=
 \int_0^1 f(x)^{2-\alpha}
       \E\bigl[f(Z)^{\alpha-1}\mid X=x\bigr]\dd x.
 \label{eq:conditional-M}
\end{equation}
In particular, for any measurable $E\subset(0,1)$, the contribution with
$X\in E$ is at most
\begin{equation}
 \|f\|_\infty^{\alpha-1}\int_E f(x)^{2-\alpha}\dd x.
 \label{eq:central-bound}
\end{equation}
\end{lemma}
\begin{proof}
The conditional density of $Z$ given $X=x$ is
$f(z)b_m(x-rz)/f(x)$. Substitution in the right side of
\eqref{eq:conditional-M} gives
$\int\!\int f(z)^\alpha b_m(x-rz)f(x)^{1-\alpha}\dd z\dd x$,
which equals \eqref{eq:master-integral}. The upper bound follows because
$f(Z)^{\alpha-1}\le\|f\|_\infty^{\alpha-1}$.
\end{proof}

\begin{lemma}[Regularity, positivity, and convergence of prefixes]
\label{lem:regularity}
For the beta family with $\theta_0,\theta_1\ge1$, $f$ is bounded,
$C^\infty$ on $\R$, positive on $(0,1)$, and zero outside $[0,1]$.
It is flat at both endpoints. Moreover, $b_m\to f$ uniformly on $\R$.
\end{lemma}
\begin{proof}
A beta density with these parameters, extended by zero, is bounded and of
bounded variation: on $(0,1)$ it is unimodal, and endpoint jumps are allowed.
Its distributional derivative is a finite signed measure. The convolution
of two bounded integrable densities is continuous. To obtain $k$ continuous
derivatives, differentiate $k$ of the first $k+2$ scaled factors in the
distributional sense; the result is a finite signed measure convolved with
a continuous compactly supported function. The remaining probability
convolution preserves continuity. This works for every $k$ and proves the
smoothness of $f$. Compact support then proves flatness.

Each finite convolution is positive on the interior of its support.
For a fixed interior point, choose a prefix long enough that the whole
possible residual displacement remains inside that support. Integrating
its positive density over the tail proves $f>0$ there. Boundedness follows
already by convolution with the first bounded factor.

Let $g$ be the convolution of the first two factors. It is bounded and
uniformly continuous. Every $b_m$ for $m\ge2$ has modulus of continuity no
larger than that of $g$. Since $0\le rZ\le r$ and $f=b_m*P_{rZ}$,
$\|f-b_m\|_\infty\le\omega_g(r)\to0$.
\end{proof}

\begin{proposition}[The finite-excess side of the transition]
\label{prop:subcritical}
For the uniform and beta channels above,
\begin{equation}
 I_1(S_m;X)=mL.
 \label{eq:Shannon}
\end{equation}
For fixed $1<\alpha\le2$,
\begin{equation}
 I_\alpha(S_m;X)-mL\longrightarrow
 \frac1{\alpha-1}\log\left[
   \left(\int_0^1 f(x)^\alpha\dd x\right)
   \left(\int_0^1 f(x)^{2-\alpha}\dd x\right)\right].
 \label{eq:subcritical-limit}
\end{equation}
At order two the second integral is $1$, the length of the support.
\end{proposition}
\begin{proof}
Boundedness and compact support give finite differential entropy. Since
$X\mid S_m=s$ has the law of $s+rZ$, the entropy difference is
$h(X)-h(rZ)=mL$, proving \eqref{eq:Shannon}.
For fixed $x\in(0,1)$, uniform convergence of $b_m$ gives
\[
 \E[f(Z)^{\alpha-1}\mid X=x]
 =\frac{\int f(z)^\alpha b_m(x-rz)\dd z}{f(x)}
 \longrightarrow \int f(z)^\alpha\dd z.
\]
Use \cref{lem:conditional}, its uniform bound on the conditional expectation,
and dominated convergence. The dominating function $f^{2-\alpha}$ is
integrable for these orders.
\end{proof}
For uniform innovations this is the previously established transition
\cite{prior}, reproduced to fix conventions and to show how the beta extension
fits it. The new issue is what happens after dominated convergence fails.
Flatness alone proves divergence for $\alpha>2$, but says nothing about its
speed. In particular, replacing the moving conditional expectation by its
pointwise limit in \eqref{eq:conditional-M} would give an infinite integral
and discard the finite-depth regularization that determines the answer.
% ed. (2026-09-30): the two predecessor theorems reproved here, by label, and the agreement of constants
\begin{ednote}
The predecessor results reproved in \cref{prop:subcritical,prop:finiteness}
are Theorem~\texttt{thm:renyi-phase} (the limit for $1<\alpha\le2$ and
divergence for $\alpha>2$) and Theorem~\texttt{thm:all-renyi} (finiteness
of every finite order at every finite depth) of
\path{../Nuclear_Bayesian_Operators_Fabius_Rvachev_Laws/article.tex}.
For uniform innovations these are second proofs of the same statements
(finiteness here by the saddle bound, there by a root-Lipschitz lemma);
the beta extension is new. The limits agree under $X_{\mathrm{old}}=2X-1$:
the product $\int f^\alpha\cdot\int f^{2-\alpha}$ is unchanged by the affine
map, and at $\alpha=2$ the predecessor's second integral is $2$, the length
of $[-1,1]$, where it is $1$ here. That article now carries a reciprocal
note at its question ``Growth above the critical R\'enyi order''.
\end{ednote}

\section{Endpoint geometry from positive simplex integrals}
\label{sec:geometry}
We work at the lower endpoint. Put
\begin{equation}
 \theta=\theta_0,\qquad \beta=\theta_1,\qquad
 H=\frac{\Gamma(\theta+\beta)}{\Gamma(\beta)},
 \qquad w_j=cq^j.
 \label{eq:endpoint-notation}
\end{equation}
At the upper endpoint the same formulas hold after interchanging
$\theta$ and $\beta$ and reflecting $x$ to $1-x$.

\subsection{The unrestricted power simplex}
For $n\ge1$, convolution of the functions
$w_j^{-\theta}u^{\theta-1}/\Bf(\theta,\beta)$ on the positive half-line gives
\begin{equation}
 A_n(t)=K_n t^{\theta n-1},\qquad
 K_n=\frac{H^n}{\Gamma(\theta n)c^{\theta n}
                         q^{\theta n(n-1)/2}},\qquad t>0.
 \label{eq:power-simplex}
\end{equation}
The identity follows by repeated application of Euler's beta integral
\cite{dlmf-beta}. It can also be proved by the change of variables
$u_j=tD_j$ on $\sum u_j=t$, where
$(D_0,\ldots,D_{n-1})$ has Dirichlet parameters $(\theta,\ldots,\theta)$.
In particular, $\E D_j=1/n$.

The actual prefix density $b_n(t)$ is $A_n(t)$ multiplied by the Dirichlet
expectation
\begin{equation}
 \E_D\left[\prod_{j<n}
 (1-tD_j/w_j)^{\beta-1}
 \boldsymbol1_{\{tD_j\le w_j\text{ for every }j\}}\right].
 \label{eq:dirichlet-factor}
\end{equation}
The integrand in \eqref{eq:dirichlet-factor} is defined to be zero whenever
any support constraint fails; the powers are evaluated only on the allowed
simplex. This avoids taking a noninteger power of a negative number.
This positive representation is useful because it avoids subtracting large
alternating truncated-power terms.

\begin{lemma}[Uniform prefix comparison in the final cell]
\label{lem:prefix}
Choose $\eta>0$ so small that $\eta q/c\le1/2$. For every $m$ and
$0<u<\eta q^m$,
\begin{equation}
 e^{-C}K_m u^{\theta m-1}\le b_m(u)\le K_m u^{\theta m-1},
 \label{eq:prefix-bounds}
\end{equation}
where $C$ depends on $q,\theta,\beta,\eta$, not on $m$ or $u$.
For uniform innovations, $\theta=\beta=1$, equality holds with $C=0$.
\end{lemma}
\begin{proof}
Here $u/w_{m-1}<\eta q/c\le1/2$, so all the indicators in
\eqref{eq:dirichlet-factor} equal one. Since $\beta\ge1$, the product is at
most one. Also $\log(1-v)\ge-2v$ for $0\le v\le1/2$, and
\[
 \sum_{j<m}\frac{uD_j}{w_j}\le\frac{u}{w_{m-1}}
                                  \le\frac{\eta q}{c}.
\]
Thus the product is bounded below by
$\exp(-2(\beta-1)\eta q/c)$, uniformly on the simplex.
\end{proof}

\subsection{A density estimate, not merely a small-ball estimate}
The denominator of the R\'enyi integrand contains a \emph{density}. A
small-ball estimate for $\Prb(X\le x)$ alone would not be an adequate
substitute. The next lemma supplies a direct density estimate to enough
precision that the $m\log m$ cancellation can be controlled.

\begin{lemma}[Endpoint logarithm with a controlled linear term]
\label{lem:density-log}
As $T\to\infty$,
\begin{equation}
 \log f(e^{-T})
 =-\frac{\theta}{2L}T^2-\frac{\theta}{L}T\log T+A_\theta T
      +O\bigl((\log(T+1))^2\bigr),
 \label{eq:density-log}
\end{equation}
where
\begin{equation}
 A_\theta=
 \frac{\theta}{L}\left(1+\log\frac{L}{\theta c}\right)
       +\frac{\log H}{L}+1-\frac{\theta}{2}.
 \label{eq:Atheta}
\end{equation}
More strongly, there is an integer-valued choice
\begin{equation}
 n(T)=\left\lfloor\frac{T+\log(T+1)}{L}\right\rfloor-K
       =\frac{T+\log T}{L}+O(1)
 \label{eq:nT}
\end{equation}
with a fixed sufficiently large integer $K$, for which
\begin{equation}
 \log f(e^{-T})=\log K_{n(T)}-(\theta n(T)-1)T+O(1).
 \label{eq:bounded-log-remainder}
\end{equation}
All constants are fixed once the innovation law and $q$ are fixed.
\end{lemma}
\begin{proof}
Write $x=e^{-T}$ and $n=n(T)$. The choice in \eqref{eq:nT} ensures that
$nq^n/x$ stays between positive finite constants. Taking $K$ sufficiently
large also makes
\begin{equation}
 \frac{x}{n}\sum_{j<n}\frac1{w_j}
 \le\frac{x}{ncw_{n-1}}\le\frac14
 \label{eq:Dirichlet-small}
\end{equation}
for all sufficiently large $T$. Increasing $K$ multiplies $w_{n-1}/x$
by a fixed geometric factor; this is compatible with $q^n/x=O(1/n)$.

For any $t\in[x-q^n,x]$, the Dirichlet expectation of
$\sum tD_j/w_j$ is at most $1/4$. Markov's inequality gives probability
at least $1/2$ that this sum is at most $1/2$. On that event every
$tD_j/w_j$ is at most $1/2$, all support restrictions are satisfied, and
\[
 \prod_{j<n}(1-tD_j/w_j)^{\beta-1}\ge e^{-(\beta-1)}.
\]
The upper bound by one remains valid everywhere. Hence
\begin{equation}
 \tfrac12 e^{-(\beta-1)}A_n(t)\le b_n(t)\le A_n(t)
 \qquad (x-q^n\le t\le x).
 \label{eq:head-density-comparison}
\end{equation}
For large $T$, $q^n/x<1/2$ and $\theta n-1>0$. Because
$nq^n/x=O(1)$,
\[
 \frac{A_n(x-q^n)}{A_n(x)}
  =(1-q^n/x)^{\theta n-1}\ge e^{-C_1}>0.
\]
The self-similar decomposition gives
$f(x)=\E[b_n(x-q^nZ)]$ with $0\le Z\le1$. Thus
$f(x)\asympc A_n(x)$ with constants independent of $T$, proving
\eqref{eq:bounded-log-remainder}.

To expand it, use
\begin{align}
 \log A_n(e^{-T})={}&-(\theta n-1)T+n\log H-\theta n\log c
       +\frac{\theta L}{2}n(n-1)-\log\Gamma(\theta n).
 \label{eq:An-log}
\end{align}
Stirling's formula \cite{dlmf-gamma}, together with
$n=(T+\log T)/L+O(1)$, gives the two leading terms in
\eqref{eq:density-log}. The coefficient of $T$ is
\[
 \frac{\theta}{L}\log L-\frac{\theta}{L}\log\theta
 -\frac{\theta}{L}\log c+\frac{\theta}{L}
 +\frac{\log H}{L}+1-\frac{\theta}{2},
\]
which is \eqref{eq:Atheta}. All remaining products contain at most two
factors of $\log(T+1)$, or are bounded; this proves the stated remainder.
\end{proof}

\begin{remark}[Why a rough log-normal equivalent is insufficient]
The statement $\log f(e^{-T})\sim-\theta T^2/(2L)$ would determine the
quadratic coefficient of $I_\alpha$, but not its linear coefficient.
The term $-\theta T\log T/L$ and the explicit $A_\theta T$ are needed:
they cancel the prefix gamma normalization and the endpoint normalization.
The bounded-remainder formula \eqref{eq:bounded-log-remainder} also leaves
room for bounded lattice-scale effects; no nonexistent constant limit is
assumed for such effects.
\end{remark}

\section{The endpoint moment and its two saddle variables}
\label{sec:saddle}
From now on $\alpha>2$, and we use
\begin{equation}
 a=\alpha-1>1,\qquad b=\alpha-2>0,
 \qquad 1+\alpha b=a^2.
 \label{eq:ab}
\end{equation}
Choose $\eta=e^{-T_0}$ sufficiently small for \cref{lem:prefix}, for the
endpoint estimates, and for the positive beta parameters below. Let
$M_{\alpha,m}^{(0)}$ be the part of \eqref{eq:master-integral} with
$0<X<\eta r$.

\subsection{The change of variables and every normalization factor}
On this region set
\begin{equation}
 x=s+rz,\qquad \lambda=\frac{s}{x},\qquad
 s=\lambda x,\quad z=(1-\lambda)x/r,
 \quad x=r e^{-W},\quad W>T_0.
 \label{eq:edge-coordinates}
\end{equation}
Here $0<\lambda<1$. The Jacobian of $(x,\lambda)\mapsto(s,z)$ is $x/r$.
Using \cref{lem:prefix}, and writing $F(T)=\log f(e^{-T})$, gives
\begin{equation}
 M_{\alpha,m}^{(0)}\asympc
 D_m\int_{T_0}^\infty e^{-(\theta m+1)W-aF(W+mL)}
 \int_0^1\lambda^{\theta m-1}
       e^{\alpha F(W-\log(1-\lambda))}\dd\lambda\dd W,
 \label{eq:edge-master}
\end{equation}
with constants independent of $m$, where
\begin{equation}
 D_m=K_m r^{\theta m},\qquad
 \log D_m=-\frac{\theta L}{2}m(m+1)+m\log H
             -\theta m\log c-\log\Gamma(\theta m).
 \label{eq:Dm}
\end{equation}
For uniform innovations the comparison in \eqref{eq:edge-master} is an
identity. The factors $x/r$ from the Jacobian and $\dd x=x\dd W$ are
responsible for the $+1$ in $\theta m+1$; omitting it changes the linear term.

The variable $W$ determines how far inside the last geometric scale the
source approaches the endpoint. The variable $\lambda$ determines how much
of this deficit is supplied by the prefix. Both must be optimized.

\subsection{Reducing the allocation integral to a beta integral}
Define the smooth comparison function
\begin{equation}
 F_0(T)=-\frac{\theta T^2}{2L}-\frac{\theta T\log T}{L}+A_\theta T.
 \label{eq:F0}
\end{equation}
By \cref{lem:density-log}, $F(T)=F_0(T)+O((\log(T+1))^2)$.

\begin{lemma}[Uniform allocation estimate]
\label{lem:allocation-integral}
Fix $0<d_1<d_2<\infty$. Uniformly for $d_1m\le W\le d_2m$,
\begin{align}
 \log\int_0^1\lambda^{\theta m-1}
 e^{\alpha F(W-\log(1-\lambda))}\dd\lambda
  =\alpha F_0(W)
    +\log\Bf\left(\theta m,\frac{\alpha\theta W}{L}+1\right)
    +O\bigl((\log(m+1))^2\bigr).
 \label{eq:allocation-beta}
\end{align}
\end{lemma}
\begin{proof}
Put $k=-\log(1-\lambda)\ge0$. For all sufficiently large $W$, the explicit
formula for $F_0$ gives
\begin{align*}
 F_0(W+k)-F_0(W)
  ={}&-\frac{\theta Wk}{L}-\frac{\theta k^2}{2L}+A_\theta k\\
 &-\frac{\theta}{L}\bigl((W+k)\log(W+k)-W\log W\bigr).
\end{align*}
The last bracket is nonnegative. The error $O(\log^2(W+k+1))$ can be bounded
by $C\log^2(W+1)+\theta k^2/(4L)+C'$. Consequently,
\begin{equation}
 F(W+k)\le F_0(W)
    -\left(\frac{\theta W}{L}-A_\theta\right)k
    +C\log^2(W+1)+C'.
 \label{eq:allocation-upper}
\end{equation}
Integration proves an upper bound with beta parameters
$\theta m$ and $\alpha\theta W/L-\alpha A_\theta+1$.
On the stated range, shifting the second beta parameter by the fixed number
$-\alpha A_\theta$ changes its logarithm by $O(1)$, by the gamma ratio
formula or Stirling's formula.

For a lower bound, restrict $\lambda$ to a fixed compact interval inside
$(0,1)$ containing the means of the beta distributions with parameters
$\theta m$ and $\alpha\theta W/L+1$, for every $W/m\in[d_1,d_2]$.
On this interval $k$ is bounded and
\[
 F(W+k)=F_0(W)-\frac{\theta Wk}{L}+O(\log^2(W+1)).
\]
The chosen interval contains at least half the beta mass for all sufficiently
large $m$, uniformly in $W$: its mean stays away from the endpoints and its
variance is $O(1/m)$, as follows from the beta integral. This proves the
matching lower bound.
\end{proof}

\subsection{Global localization before local expansion}
It is not enough to find a favorable stationary point. The very small
endpoint tail and the regions where $\lambda$ approaches zero or one must
also be controlled.

\begin{lemma}[Localization of the radial integral]
\label{lem:localization}
There are constants $0<d_1<bL<d_2<\infty$ such that the part of
\eqref{eq:edge-master} with $W\notin[d_1m,d_2m]$ is exponentially smaller,
at speed $m^2$, than its whole value. On the retained interval its logarithm
differs by $O(\log^2(m+1))$ from the logarithm of
\begin{equation}
 \int_{d_1m}^{d_2m}\exp(E_m(W))\dd W,
 \label{eq:Em-integral}
\end{equation}
where
\begin{align}
 E_m(W)={}&\log D_m-(\theta m+1)W
       +\alpha F_0(W)-aF_0(W+mL)\nonumber\\
       &+\log\Bf\left(\theta m,\frac{\alpha\theta W}{L}+1\right).
 \label{eq:Em}
\end{align}
\end{lemma}
\begin{proof}
Take $T_0$ large enough that the second beta parameter in the upper bound
\eqref{eq:allocation-upper} is at least one. That beta integral is at most
$1/(\theta m)$. The quadratic part of the remaining exponent, including the
quadratic term of $\log D_m$, is exactly
\begin{equation}
 \frac{\theta L}{2}abm^2
             -\frac{\theta}{2L}(W-bmL)^2.
 \label{eq:quadratic-landscape}
\end{equation}
All other terms are bounded in absolute value by
$C(m+W+1)\log(m+W+2)$. For $W$ of order $m$, this is $o(m^2)$.
For $W$ much larger than $m$, it is dominated by a fixed positive fraction of
$W^2$, uniformly once $m$ is large. Thus choosing $d_1<bL<d_2$ away from
$bL$ and $d_2$ sufficiently large gives a Gaussian-integrable upper bound
outside the interval, with a loss $c_1m^2$ from the maximum in
\eqref{eq:quadratic-landscape}.

A lower bound of size
$\exp(\theta Labm^2/2-O(m\log(m+1)))$ is obtained by taking $W$ in a fixed
width interval about $bmL$ and $\lambda$ in a fixed compact subinterval of
$(0,1)$. This proves that the omitted tails are relatively
$\exp(-c_2m^2)$ after increasing $m$ if necessary. Finally,
\cref{lem:allocation-integral} and the uniform bound
$F(W+mL)-F_0(W+mL)=O(\log^2(m+1))$ prove the assertion on the retained
interval. Multiplicative constants from \cref{lem:prefix} contribute only
$O(1)$ to its logarithm.
\end{proof}

\subsection{Evaluation of the radial saddle}
\begin{lemma}[Saddle evaluation at logarithmic-square accuracy]
\label{lem:radial}
With $W_0=bmL$,
\begin{equation}
 \log M_{\alpha,m}^{(0)}=E_m(W_0)+O(\log^2(m+1)).
 \label{eq:radial-evaluation}
\end{equation}
\end{lemma}
\begin{proof}
Differentiate the explicit comparison exponent \eqref{eq:Em}, not the
unknown error in $F$. Since
$F_0''(T)=-\theta/L-\theta/(LT)$, and the second derivative of the log beta
term is $O(1/m)$ when both parameters are proportional to $m$, one has
\begin{equation}
 E_m''(W)=-\frac{\theta}{L}+O(1/m)
 \label{eq:Em-concavity}
\end{equation}
uniformly on the interval in \cref{lem:localization}. The beta derivative
assertion follows directly from
$\Bf(u,v)=\Gamma(u)\Gamma(v)/\Gamma(u+v)$ and the usual differentiated
Stirling estimates. Alternatively, the trigamma series gives it directly.
The same formulas give
\begin{equation}
 E_m'(W_0)=O(\log(m+1)).
 \label{eq:Em-first-derivative}
\end{equation}
In fact its leading term is $-\theta\log m/L$, but its value is not needed.
The maximum of $E_m$ is therefore within $O(\log(m+1))$ of $W_0$ and its
height differs from $E_m(W_0)$ by $O(\log^2(m+1))$.
Uniform strict concavity bounds the integral above by a constant times the
exponential of that height. Integrating over an interval of length one about
$W_0$ supplies a lower bound with logarithmic loss at most $O(\log(m+1))$.
Apply \cref{lem:localization}.
\end{proof}

\subsection{The cancellation that produces universality}
The next evaluation is the algebraic center of the proof. All terms of size
$m\log m$ must be retained before cancellation.

\begin{proposition}[One-endpoint asymptotic]
\label{prop:endpoint-asymptotic}
Let $K(\alpha)=\alpha b\log\alpha-a^2\log a=aB(\alpha)$. Then
\begin{align}
 \log M_{\alpha,m}^{(0)}
  & =\frac{\theta L}{2}abm^2
       +\bigl(-aL+\theta K(\alpha)\bigr)m
       +O(\log^2(m+1)),\label{eq:endpoint-M}\\
 \log J_{\alpha,m}^{(0)}
  & =\frac{\theta L}{2}abm^2
       +\theta K(\alpha)m+O(\log^2(m+1)).
 \label{eq:endpoint-J}
\end{align}
Here $J^{(0)}=r^{-a}M^{(0)}$ is the corresponding contribution to the
unnormalized R\'enyi moment.
\end{proposition}
\begin{proof}
By \cref{lem:radial}, it suffices to evaluate $E_m(bmL)$. Stirling's formula
and $1+\alpha b=a^2$ give
\begin{equation}
 \log\Bf(\theta m,\alpha\theta bm+1)
 =\theta m\bigl[\alpha b\log(\alpha b)-2a^2\log a\bigr]
       +O(\log(m+1)).
 \label{eq:beta-stirling}
\end{equation}
The quadratic coefficient is
\[
 \theta L\left(-\frac12-b-\frac{\alpha b^2}{2}
                                    +\frac{a^3}{2}\right)
       =\frac{\theta L}{2}ab.
\]
The $m\log m$ coefficient is $-\theta+\theta(a^2-\alpha b)=0$.
Before simplification, the remaining coefficient of $m$ is
\begin{align}
 C={}&-\frac{\theta L}{2}+\log H-\theta\log c
          -\theta\log\theta+\theta-bL-A_\theta L\nonumber\\
    &+\theta\bigl[-\alpha b\log(bL)+a^2\log(aL)\bigr]
      +\theta\bigl[\alpha b\log(\alpha b)-2a^2\log a\bigr].
 \label{eq:raw-linear}
\end{align}
Inserting \eqref{eq:Atheta} cancels $\log H$, $\log c$, $\log\theta$, and
$\log L$. The result is
\begin{equation}
 C=-aL+\theta\bigl[\alpha b\log\alpha-a^2\log a\bigr].
 \label{eq:linear-cancelled}
\end{equation}
This proves \eqref{eq:endpoint-M}. Multiplying by $r^{-a}$ adds $aLm$ to the
logarithm and proves \eqref{eq:endpoint-J}.
\end{proof}
The parameter $H$ contains the normalization at the chosen innovation
endpoint and depends on the opposite beta exponent. Its cancellation is
why the linear coefficient is not contaminated by that normalization.
Only the endpoint exponent $\theta$ survives.

\section{Assembling the two endpoints and the interior}
\label{sec:assembly}
Let $M^{(1)}$ be the contribution with $1-\eta r<X<1$, and let
$M^{\mathrm{mid}}$ be the remaining contribution. Reflection changes the
lower-endpoint exponent $\theta_0$ to $\theta_1$, so
\cref{prop:endpoint-asymptotic} applies to $M^{(1)}$ as well.

\begin{lemma}[The interior is negligible above order two]
\label{lem:middle}
For fixed $\alpha>2$,
\begin{equation}
 \log M_{\alpha,m}^{\mathrm{mid}}
 \le\frac{\theta_*L}{2}b m^2+O(m\log(m+1)).
 \label{eq:middle-bound}
\end{equation}
\end{lemma}
\begin{proof}
On a fixed closed subinterval of $(0,1)$, $f$ is bounded below by a positive
constant. On the two remaining portions of
$[\eta r,1-\eta r]$, \cref{lem:density-log} and its reflected version give
\[
 -\log f(x)\le\frac{\theta_*L}{2}m^2+O(m\log(m+1)),
\]
since the corresponding logarithmic endpoint distance is at most
$mL+T_0$. The estimate \eqref{eq:central-bound} now proves the claim.
\end{proof}
The largest endpoint coefficient in \eqref{eq:endpoint-M} is
$\theta_*Lab/2$, which exceeds the middle coefficient
$\theta_*Lb/2$ by $\theta_*Lb^2/2>0$.
Hence the middle contribution is relatively exponentially small at speed
$m^2$. Taking the logarithm of the sum of the two endpoint terms selects
% ed. (2026-09-30): "the larger $\theta_e$" read "the endpoint with the larger exponent"; $\theta_e$ is defined only in Section 6.
the endpoint with the larger exponent $\theta_0$ or $\theta_1$; when they are equal, summing them changes the logarithm
only by $\log2$. Finally divide \eqref{eq:endpoint-J} by $a$.
This proves \cref{thm:beta-main}, and \cref{thm:uniform} is its
$\theta_0=\theta_1=1$ case.

\begin{corollary}[Quantitative endpoint selection]
\label{cor:selection}
The two endpoint contributions satisfy
\begin{equation}
 \log\frac{M_{\alpha,m}^{(0)}}{M_{\alpha,m}^{(1)}}
 = (\theta_0-\theta_1)
    \left[\frac{Lab}{2}m^2+K(\alpha)m\right]
       +O(\log^2(m+1)).
 \label{eq:endpoint-ratio}
\end{equation}
The normalized middle contribution satisfies
\begin{equation}
 \limsup_{m\to\infty}\frac1{m^2}
 \log\frac{M_{\alpha,m}^{\mathrm{mid}}}{M_{\alpha,m}}
 \le-\frac{\theta_*L}{2}b^2.
 \label{eq:middle-probability}
\end{equation}
If the endpoint exponents are equal, each endpoint carries asymptotically
one half of the tilted mass. If they differ, the larger-exponent endpoint
carries asymptotically all of it.
\end{corollary}

\begin{proposition}[Every finite order at every finite depth]
\label{prop:finiteness}
For each fixed $m\ge1$ and every finite $\alpha>1$, $I_\alpha(S_m;X)$ is
finite for the beta channels considered here.
\end{proposition}
\begin{proof}
The middle contribution is finite because its denominator is bounded below.
For an endpoint, use \eqref{eq:edge-master} and the bound
\eqref{eq:allocation-upper}. With $m$ fixed, the coefficient of $W^2$ in
$\alpha F_0(W)-aF_0(W+mL)$ is $-\theta/(2L)$. All remaining terms are
$O(W\log(W+1))$ or linear in $W$. The radial integral is therefore finite.
The inner allocation integral is bounded by a beta integral once $T_0$ is
large enough; the omitted finite $W$ interval can be treated by increasing
$T_0$ and regarding it as part of the middle region. Reflect for the other
endpoint.
\end{proof}
Thus quadratic divergence with depth is fully compatible with existence
at every fixed depth. No finite-$m$ singularity at $\alpha=2$ is asserted.

\section{Which rare configurations create the large moment?}
\label{sec:tilt}
The saddle variables also have a probabilistic interpretation. Define a
probability measure $\mathsf Q_{\alpha,m}$ on $(s,z)$ by
\begin{equation}
 \mathsf Q_{\alpha,m}(\dd s\dd z)
 =\frac{b_m(s)f(z)^\alpha f(s+rz)^{-a}}{M_{\alpha,m}}
                          \dd s\dd z.
 \label{eq:tilted-measure}
\end{equation}
This is the \emph{R\'enyi-tilted} measure, not the ordinary joint law of the
prefix and the tail. Statements about its typical configurations describe
which events dominate a likelihood moment; they do not say that an ordinary
sample of $X$ usually lies near an endpoint.

For endpoint $e=0$, let
$d=x$, $u=s$, and $v=z$. For endpoint $e=1$, let
$d=1-x$, $u=(1-r)-s$, and $v=1-z$. In either case
\begin{equation}
 d=u+rv,\qquad \lambda=u/d,\qquad W=\log(r/d),
 \label{eq:deficit-coordinates}
\end{equation}
and we condition on the endpoint cell $d<\eta r$. Write
$\mathsf Q_{\alpha,m}^{(e)}$ for this conditional measure and
$\theta_e$ for its innovation exponent.

\begin{theorem}[Radial large deviations inside an endpoint cell]
\label{thm:radial-LDP}
Under $\mathsf Q_{\alpha,m}^{(e)}$, the variables
$W/(mL)$ satisfy a large-deviation principle on $[0,\infty)$ with speed
$m^2$ and good rate function
\begin{equation}
 R_e(t)=\frac{\theta_eL}{2}(t-b)^2,
 \qquad t\ge0.
 \label{eq:radial-rate}
\end{equation}
In particular $W/(mL)\to b$ in probability.
\end{theorem}
\begin{proof}
On every compact subinterval of $(0,\infty)$, put $W=mLt$ in the
endpoint integral. \Cref{lem:density-log,lem:prefix} and the beta integral
estimate show that its logarithmic density, divided by $m^2$, converges
uniformly to
\[
 \frac{\theta_eL}{2}\bigl(ab-(t-b)^2\bigr).
\]
The factor $mL$ from the change of variables is negligible on this scale.
The upper estimates in \cref{lem:localization} give exponential tightness
at $+\infty$. At $t=0$, a lower bound is obtained by integrating over a
fixed positive-length interval of $W>T_0$ and a compact interval of
$\lambda$; its coefficient is $\theta_eLb/2$, the same as the continuous
extension of the displayed expression at zero. Upper bounds near zero
follow from the same quadratic bound with errors $O(m\log(m+1))$.

Cover a compact closed set by finitely many small intervals and apply the
uniform upper bound; integrate on a small interval about any point of an
open set for the lower bound. Exponential tightness handles the unbounded
remainder. Subtract the normalizing coefficient $\theta_eLab/2$ from
\eqref{eq:endpoint-M}. This gives the upper and lower large-deviation
bounds with \eqref{eq:radial-rate}. Its sublevel sets are compact.
\end{proof}

\begin{theorem}[Universal allocation and its exponential penalty]
\label{thm:allocation-LDP}
Under $\mathsf Q_{\alpha,m}^{(e)}$, the variables $\lambda$ satisfy a
large-deviation principle with speed $m$ and rate function
\begin{equation}
 A_e(\lambda)=\theta_e\left[
   \log\frac{\lambda_*}{\lambda}
   +\alpha b\log\frac{1-\lambda_*}{1-\lambda}\right],
 \quad 0<\lambda<1,\qquad \lambda_*=\frac1{a^2},
 \label{eq:allocation-rate}
\end{equation}
with value $+\infty$ at zero and one. In particular,
\begin{equation}
 \frac{u}{d}\longrightarrow\frac1{(\alpha-1)^2},\qquad
 \frac{rv}{d}\longrightarrow\frac{\alpha(\alpha-2)}{(\alpha-1)^2}
 \label{eq:allocation-limits}
\end{equation}
in probability. The rate can equivalently be written as
$\theta_ea^2\KL(\operatorname{Bern}(\lambda_*)\Vert
\operatorname{Bern}(\lambda))$.
\end{theorem}
\begin{proof}
For $\lambda$ in any fixed compact subinterval of $(0,1)$, $k=-\log(1-\lambda)$
is bounded. Use $F(W+k)=F_0(W)-\theta_eWk/L+O(\log^2(m+1))$ uniformly
for $W$ proportional to $m$. The proof of \cref{lem:radial} works with
$\lambda$ fixed in that compact interval: the second derivative of its
comparison exponent is $-\theta_e/L+O(1/m)$, and its first derivative at
$bmL$ is $O(\log(m+1))$ uniformly there. After integrating in $W$, the
$\lambda$-dependent part of the logarithm is therefore
\begin{equation}
 \theta_e m\bigl[\log\lambda+\alpha b\log(1-\lambda)\bigr]
                           +O(\log^2(m+1)),
 \label{eq:allocation-logdensity}
\end{equation}
up to a term independent of $\lambda$.
This strictly concave bracket has its unique maximum at
$\lambda_*=1/(1+\alpha b)=a^{-2}$. Subtracting that maximum gives
\eqref{eq:allocation-rate}.

For completeness, the endpoints cannot hide additional mass. First remove
$|W/(mL)-b|>\delta$ using \cref{thm:radial-LDP}; this costs an amount
exponentially small at speed $m^2$. On the remainder,
\eqref{eq:allocation-upper} bounds the allocation integral by beta powers
with the second exponent at least $\alpha\theta_e(b-\delta)m-O(1)$.
The integral over $\lambda<\eps$ is bounded by a constant times
$\eps^{\theta_e m}$ times the common radial factor; the integral over
$\lambda>1-\eps$ has a bound with a positive multiple of
$m\log\eps$ from the second beta exponent. To see that the common radial factor costs only $O(m)$ relative to the
normalization, drop the beta term from \eqref{eq:Em} and apply the same
strict-concavity argument as in \cref{lem:radial}. At $W=bmL$ its
$m\log m$ coefficient is still $-\theta_e+\theta_e(a^2-\alpha b)=0$;
its logarithmic integral is therefore
$\theta_eLabm^2/2+O(m)$. Equation~\eqref{eq:endpoint-M} gives the same
quadratic normalization. Their logarithmic ratio is thus $O(m)$,
independently of $\eps$. Letting $\eps\downarrow0$ makes both upper bounds arbitrarily
negative at speed $m$. Thus the laws are exponentially tight away from
zero and one at that speed. Compact-set upper and open-set lower bounds
now follow by the same interval-covering argument applied to
\eqref{eq:allocation-logdensity}. The two probability limits follow from
uniqueness of the minimum, and the Bernoulli expression is an algebraic
rewriting of the rate.
\end{proof}

\begin{corollary}[The source and tail scales]
\label{cor:scales}
Conditional on either endpoint cell, under the tilted measure,
\begin{equation}
 \frac{-\log d}{mL}\to\alpha-1,
 \qquad \frac{-\log v}{mL}\to\alpha-2.
 \label{eq:scales}
\end{equation}
Equivalently,
$d=q^{(\alpha-1)m+o_{\mathsf Q}(m)}$ and
$v=q^{(\alpha-2)m+o_{\mathsf Q}(m)}$.
\end{corollary}
\begin{proof}
The first follows from $-\log d=mL+W$. For the second, use
$v=(1-\lambda)e^{-W}$ and \cref{thm:allocation-LDP}; the logarithm of
$1-\lambda$ is bounded in probability and is therefore negligible after
division by $mL$.
\end{proof}
For example, at $\alpha=3$ the dominant prefix fraction is $1/4$, and the
residual tail contributes $3/4$ of the endpoint deficit. The source lies
on scale $q^{2m+o(m)}$ and the rescaled tail on scale $q^{m+o(m)}$.
The allocation fractions are independent of $q$, the beta exponents, and
which endpoint wins.

\section{Consequences, checks, and the limits of the expansion}
\label{sec:consequences}

\subsection{The linear coefficient is positive and universal}
\begin{proposition}
$B(\alpha)$ is strictly increasing on $(2,\infty)$, with
\begin{equation}
 B(2+)=0,\qquad \lim_{\alpha\to\infty}B(\alpha)=1.
 \label{eq:B-limits}
\end{equation}
In particular $0<B(\alpha)<1$. Near the critical order,
\begin{equation}
 B(2+\eps)=(2\log2-1)\eps+O(\eps^2).
 \label{eq:B-near2}
\end{equation}
\end{proposition}
\begin{proof}
Put $a=\alpha-1>1$. Then
$B=(a-a^{-1})\log(a+1)-a\log a$, so
\[
 a^2B'(a)=a^2\log(1+1/a)-a+\log(a+1).
\]
The inequality $\log(1+t)\ge t-t^2/2$ for $0\le t\le1$ gives
$a^2B'(a)\ge\log(a+1)-1/2>0$. The endpoint limits and the expansion at
$a=1$ follow by direct Taylor expansion.
\end{proof}
For $q=1/2$, $\alpha=3$ and uniform innovations,
\begin{equation}
 I_3(S_m;X)=\frac{\log2}{2}m^2
       +\left(\frac32\log3-2\log2\right)m+O(\log^2(m+1)).
 \label{eq:alpha3}
\end{equation}
The linear coefficient is approximately $0.2616240719$. For
$\operatorname{Beta}(2,5)$ innovations the two coefficients are both
multiplied by five, and the upper endpoint dominates. This is an application
of the theorem, not a numerical fit.

\subsection{Why neither endpoint flatness nor a single saddle is enough}
Endpoint flatness explains divergence of negative powers of $f$, but all
finite-depth R\'enyi integrals remain finite. The finite-prefix factor
$u^{\theta m-1}/\Gamma(\theta m)$ regulates the apparent singularity.
The tail factor contributes another logarithmic-square penalty. Their
competition yields \eqref{eq:quadratic-landscape}.

Optimizing only the radial exponent finds the $m^2$ coefficient, but loses
the beta allocation entropy, which is of order $m$. Replacing
$\lambda$ by an arbitrary constant gives a strictly smaller linear term
unless it equals $a^{-2}$. Conversely, differentiating a small-tail
asymptotic with an uncontrolled error is unnecessary: the proof differentiates
only $F_0$ and the beta comparison exponent. This is why a mere value estimate
with a quantified remainder is sufficient.

\subsection{Nonuniform limits}
At $\alpha=2$, \cref{prop:subcritical} gives
$I_2=mL+\log\int f^2+o(1)$. At every fixed $\alpha>2$, the leading order is
$m^2$. There is no contradiction: the constants and localization intervals
in the supercritical proof depend on $b=\alpha-2$. They degenerate as
$b\downarrow0$.

In particular, \eqref{eq:beta-main} does not solve the joint scaling
$\alpha=2+u/m$. It also is not uniform as $q\uparrow1$, as
$\alpha\to\infty$, or as beta parameters vary. Nor does it give an exact
finite-depth formula for $I_\alpha$. A rigorous asymptotic statement with a
specified range should not be advertised as any of these stronger claims.

\section{Reproducible algebra and tail-density diagnostics}
\label{sec:computations}
The accompanying \path{checks/verify.py} separates exact arithmetic from
floating-point diagnostics. It requires only the packages in
\path{requirements.txt}. No plot or decimal calculation is used as a proof
of an analytic assertion.

\subsection{Exact checks of the cancellation and the allocation}
SymPy checks the quadratic identity in \cref{prop:endpoint-asymptotic},
the complete cancellation from \eqref{eq:raw-linear} to
\eqref{eq:linear-cancelled}, and stationarity of
$\log\lambda+\alpha(\alpha-2)\log(1-\lambda)$ at
$\lambda=(\alpha-1)^{-2}$. The symbolic endpoint constants $c,H,\theta,L$
are retained independently until the cancellation step.

For the dyadic uniform law, moments $\mu_n=\E X^n$ are calculated exactly
as rational numbers using
\begin{equation}
 \mu_0=1,\qquad
 \mu_n=\frac1{2^n-1}\sum_{j=0}^{n-1}
             \binom nj\frac{\mu_j}{n-j+1}.
 \label{eq:moment-recurrence}
\end{equation}
This follows by expanding $X\overset d=(U+X')/2$.
The checks include $\E X=1/2$ and $\Var X=1/36$.

\subsection{Exact dyadic densities}
\begin{proposition}[A finite exact evaluation identity]
\label{prop:dyadic}
For integers $n\ge1$ and $1\le k\le2^n$, the density of the dyadic uniform
series satisfies
\begin{equation}
 f(k2^{-n})=
 \frac{2^{(-n^2+3n)/2}}{(n-1)!}
 \sum_{j=0}^{k-1}(-1)^{s_2(j)}\E[(k-j-X)^{n-1}],
 \label{eq:dyadic-density}
\end{equation}
where $s_2(j)$ is the sum of the binary digits of $j$.
For $k=1$ this simplifies to
\begin{equation}
 f(2^{-n})=\frac{2^{(-n^2+3n)/2}}{(n-1)!}\mu_{n-1}.
 \label{eq:dyadic-one}
\end{equation}
\end{proposition}
\begin{proof}
Inclusion--exclusion on the boxes of the first $n$ uniforms gives
\[
 b_n(t)=\frac{1}{(n-1)!\prod_{j=1}^n2^{-j}}
       \sum_{j=0}^{2^n-1}(-1)^{s_2(j)}(t-j2^{-n})_+^{n-1}.
\]
Convolve with $2^{-n}X'$. At $t=k2^{-n}$, exactly the terms $j<k$ can
contribute; in each of them $k-j-X'\ge0$. Factoring out
$2^{-n(n-1)}$ gives \eqref{eq:dyadic-density}. Symmetry $1-X\overset d=X$
gives \eqref{eq:dyadic-one}. For $n=1$ the endpoint values of the
piecewise-constant prefix are irrelevant to the convolution.
\end{proof}
This is a convenient verification identity, not a claim to a new
arithmetic theory of dyadic Fabius values; that theory is classical
\cite{arias}. The script checks $f(1/2)=2$, $f(1/4)=1$, $f(1/8)=5/36$,
and 24 independent dyadic refinement equalities.
% ed. (2026-09-30): uncited Lean counterparts of the dyadic values (names checked in Analysis/FabiusFunction/Lean)
\begin{ednote}
At $q=1/2$ these values are machine-checked in ProveIt's Lean development
(\path{Analysis/FabiusFunction/Lean/FabiusFunction/}), which the article
does not cite. There $f(x)=2\,\mathrm{up}(2x-1)$
(\path{Fabius.ProbabilityRepresentation.geometricUniformDensity_one_half_eq_rvachevUp}),
equivalently $f(x)=2F(2x)$ on $[0,1/2]$ for the Fabius distribution
function $F$; the exact dyadic values of $F$ are computed by a verified
evaluator (\path{Fabius.fabiusDyadicValue_cast}), for instance
$F(1/4)=5/72$ and $F(3/4)=67/72$ (\path{Fabius.fabiusReal_quarter},
\path{Fabius.fabiusReal_three_quarters}), giving $f(1/8)=5/36$ and
$f(3/8)=67/36$; and $f(1/2)=2$ is \path{Fabius.deriv_fabiusReal_half}.
No R\'enyi statement of the article has a Lean counterpart.
\end{ednote}

\subsection{A tail evaluator with an analytic truncation bound}
For a non-dyadic $x\in(0,1/2]$, choose $n$ and put $k=2^nx$,
$K=\lfloor k\rfloor$. In the same inclusion--exclusion formula,
all terms $j<K$ are exactly computable from ordinary moments. There is
only one additional possible term, with $j=K$, whose absolute contribution
before the common prefactor is at most
\begin{equation}
 (k-K)^{n-1}.
 \label{eq:density-remainder}
\end{equation}
All terms with $j>K$ vanish. Thus if $S$ is the computed signed moment sum
and $S>(k-K)^{n-1}$, the relative omitted-term error is at most
$(k-K)^{n-1}/(S-(k-K)^{n-1})$. This is an analytic truncation estimate.
The floating-point implementation does not add interval bounds for rounding,
so its decimal output is not an end-to-end certified enclosure.

\subsection{Executed point diagnostics}
For $\alpha=3$, $q=1/2$, take $h=\lfloor\log_2m\rfloor$,
$x=2^h2^{-2m}$ and $\lambda=1/4$. Then
$z=3\,2^{-(m-h+2)}$, so both densities in the transformed endpoint integrand
are exactly rational by \eqref{eq:dyadic-density}. Table~\ref{tab:saddle}
compares its logarithmic point value, converted to the information scale,
with the two-term expression in \eqref{eq:alpha3}.
\begin{table}[htbp]
\centering\small
\caption{Exact dyadic point diagnostics, with logarithms rounded for display.
The point value is \emph{not} the integrated R\'enyi information.}
\label{tab:saddle}
\begin{tabular}{@{}rrrr@{}}\toprule
$m$&Point value&Two-term expression&Difference\\\midrule
\inputtable checks/saddle_rows.tex 
\bottomrule\end{tabular}
\end{table}
The modest residuals are consistent with the logarithmic-square error,
but finitely many such values cannot establish an asymptotic theorem.

\subsection{Windowed quadrature diagnostics}
For the same $q,\alpha$, the optional quadrature evaluates both endpoint
cells using Gauss--Legendre nodes in $W$ and $\lambda$, with
\[
 W\in[\log4,\,m\log2-\log m+12].
\] It excludes the middle region and truncates the
radial tail. The output explicitly labels these computations as
\emph{not certified full information values}. The JSON files record the
number of nodes, precision, integration window, and omitted-density-term
bound. Comparing node counts is a numerical stability check, not a substitute
for the proved global estimates.

To reproduce the exact checks and point table, run
\begin{verbatim}
python -m pip install -r requirements.txt
python checks/verify.py --dps 100
\end{verbatim}
An optional numerical run is
\begin{verbatim}
python checks/verify.py --quadrature 8 12 --nodes 24 --dps 80
\end{verbatim}
The PDF can be rebuilt independently of Python, because the generated table
is included:
\begin{verbatim}
pdflatex -interaction=nonstopmode -halt-on-error article.tex
pdflatex -interaction=nonstopmode -halt-on-error article.tex
\end{verbatim}

\section{Proof audit and a formalization route}
\label{sec:audit}
\begin{longtable}{@{}>{\raggedright\arraybackslash}p{.25\textwidth}
 >{\raggedright\arraybackslash}p{.45\textwidth}
 >{\raggedright\arraybackslash}p{.24\textwidth}@{}}
\toprule
Claim&Dependency and proof boundary&Status\\\midrule
\endhead
Uniform two-term law&Special case of the beta theorem; endpoints plus the
middle estimate&Proved in this article\\
Beta universality&Positive simplex density comparison, controlled endpoint
logarithm, beta allocation, radial comparison saddle&Proved for
$\theta_0,\theta_1\ge1$ only\\
Endpoint selection&Comparison of the two separately computed endpoint
moments&Proved\\
Radial and allocation laws&The same integral with restricted radial or
allocation domains; exponential tightness&Proved for the tilted measures\\
Finite-order existence&Quadratic decay of the logarithmic radial integrand
at fixed $m$&Proved\\
Critical and subcritical excess&Conditional representation and dominated
convergence; uniform case already in the predecessor&Reproved and extended
to the beta family\\
Exact algebra and dyadic values&Fraction arithmetic and symbolic identities&Executed
finite checks\\
Numerical endpoint quadrature&Floating-point density evaluation and a finite
integration window&Diagnostic only\\
Lean or Rocq verification&No theorem-prover files are included&Not performed\\
Joint critical crossover&Requires uniform estimates as $\alpha\downarrow2$&Open here\\
\bottomrule
\end{longtable}

A formalization can be staged without treating the whole asymptotic argument
as an indivisible task. The first layer comprises the independent-prefix
identity, density changes of variables, \eqref{eq:master-integral}, and
\eqref{eq:conditional-M}. A second layer formalizes the positive simplex
integral and its Dirichlet expectation representation. The endpoint lemma
then becomes a quantitative inequality about nonnegative integrals, Markov's
inequality, and gamma asymptotics. Finally, a reusable one-dimensional
comparison-saddle lemma can package strict concavity with bounded value
errors. The proof carefully avoids differentiating an uncontrolled remainder,
which also simplifies such an interface.

Exact symbolic cancellations are particularly suitable early targets:
\eqref{eq:raw-linear} contains all normalizations and can be checked
independently of probability theory. Numerical quadrature should remain
outside the trusted theorem boundary unless an interval integration checker
is added.

\section{Further research questions and proposed routes}
\label{sec:future}

\begin{problem}[The logarithmic term and lattice-scale modulation]
Improve the $O(\log^2m)$ error to an explicit logarithmic correction and a
bounded remainder, or identify the actual next scale. The stronger endpoint
representation \eqref{eq:bounded-log-remainder} is a starting point.
A refined density expansion may contain bounded modulation from the
geometric lattice; it should be retained rather than silently replaced by
a constant.
\end{problem}

\begin{problem}[A uniform critical crossover]
Determine the asymptotics for $\alpha=2+u/m$, and identify the correct
transition scale and profile. The existing fixed-order theorem is not
uniform there: the saddle $W\sim(\alpha-2)mL$ leaves the range in which
both beta parameters are proportional to $m$, and the allocation approaches
an endpoint. A uniform analysis must match the ordinary bulk contribution,
the final-cell contribution, and the allocation boundary layer.
\end{problem}
% ed. (2026-09-30): the same question is open in the predecessor
\begin{ednote}
This is the predecessor's question ``Critical crossover'' (the scaling
$\alpha=2+u/m$), in
\path{../Nuclear_Bayesian_Operators_Fabius_Rvachev_Laws/article.tex};
it remains open there and here.
\end{ednote}

\begin{problem}[Fluctuations beyond large deviations]
Find the limiting distribution of the centered radial variable and of
$\sqrt m(\lambda-\lambda_*)$. The allocation entropy suggests Gaussian
behavior in the latter variable, but the bounded geometric-lattice
modulation in the radial density may survive centering. Prove joint limits
rather than inferring them from the leading saddle alone.
\end{problem}

\begin{problem}[Endpoint exponents below one]
Extend the beta theorem to $0<\theta_0,\theta_1<1$ or mixed parameter
ranges. The final-cell comparison still has a natural form, but the
unrestricted power-simplex upper bound in \cref{lem:density-log} no longer
follows from $(1-u)^{\beta-1}\le1$. A positive tilted-Laplace argument or a
carefully weighted simplex estimate is required. No extension to that
range has been assumed in this article.
\end{problem}

\begin{problem}[Regularly varying endpoint masks]
Replace beta innovations by densities
$h(u)\sim d_0u^{\theta_0-1}\ell_0(u)$ and
$h(1-u)\sim d_1u^{\theta_1-1}\ell_1(u)$, with slowly varying factors.
Determine exactly when the quadratic law persists, whether $B(\alpha)$
remains universal, and which slowly varying factors create an additional
$m\log\log m$ or other term. The cancellation of endpoint constants above
does not by itself imply cancellation of varying factors.
\end{problem}

\begin{problem}[Nongeometric weights]
For $X=\sum w_jU_j$ with $\log(1/w_j)\sim Lj^\rho$, determine the depth
scale and endpoint-selection law. Geometric weights make the small-tail
cost quadratic in a logarithmic coordinate and the prefix normalization
quadratic in $m$. A general weight profile should lead to a different
variational problem and perhaps a different power of $m$.
\end{problem}

\begin{problem}[Certified finite-depth R\'enyi computation]
Combine the one-term density enclosure in \eqref{eq:density-remainder}
with interval arithmetic and explicit domain truncation bounds. The goal
is a certified upper and lower enclosure for the full $I_\alpha$, not
merely a stable endpoint quadrature. Such a tool could also test the
critical crossover before its proof is complete.
\end{problem}

\begin{problem}[Other R\'enyi information definitions]
Analyze Sibson and Arimoto information for the same channel. Their
optimization or different marginal weighting changes the endpoint balance.
The order-two transition and the coefficient in \eqref{eq:beta-main}
should not be transferred to those definitions without a separate proof.
\end{problem}

\begin{problem}[Products, common digits, and dependent innovations]
For vector-valued series, compare independent-coordinate products with
common-digit geometric convolutions. In a product channel, ordinary
additivity gives a direct benchmark; for common digits the relevant
geometry is a boundary face or cone rather than two scalar endpoints.
Determine how face dimension and local density exponents replace the
scalar selection rule.
\end{problem}

\section{Provenance and conclusion}
\label{sec:provenance}
The source target was read through the GitHub connector at the pinned
repository commit
\begin{center}\small
\texttt{a4268e78ebd0f6bf71ba609c07f4b3415748f532}.
\end{center}
The inspected predecessor is
\begin{quote}\small\path{Analysis/FabiusFunction/docs/}\allowbreak
\path{semi-formalized-research-frontiers/drafts/}\allowbreak
\path{inverse-and-sampling/}\allowbreak
\path{Nuclear_Bayesian_Operators_Fabius_Rvachev_Laws/article.tex}.
\end{quote}
Its returned blob SHA is
\begin{center}\small
\texttt{960890e910fcb8bfb857dfb0e08b4303c3f139f5}.
\end{center}
This records a specific inspected source, not a claim that the repository
will remain unchanged. The current article is a separate deliverable; no
repository files were modified.

The principal advance is sharper than the initially proposed quadratic
upper and lower bounds. The fixed-supercritical information has two explicit
terms, no $m\log m$ term, and a logarithmic-square remainder. Its second
coefficient is independent of the contraction parameter. The beta extension
shows that this is an endpoint phenomenon rather than a peculiarity of one
special function, while the tilted allocation law identifies the rare
configurations that produce it. The unresolved crossover and finer lattice
terms are separate, precisely stated next questions.

\begin{thebibliography}{9}
\bibitem{repo}
V.~Reshetnikov, \emph{ProveIt}, GitHub repository, inspected 30 September 2026.
\url{https://github.com/VladimirReshetnikov/ProveIt}.

\bibitem{prior}
\emph{Beyond Polynomial Diagonals: Nuclear Bayesian Operators, Exact Null
Modes, and a R\'enyi Phase Transition for Fabius--Rvachev Laws}, research
article prepared for Vladimir Reshetnikov, 29 September 2026, with repository
editorial notes. Pinned source and blob recorded in \cref{sec:provenance}.
In particular, the further question ``Growth above the critical R\'enyi
order.'' Available within the repository \cite{repo} at the recorded path.

\bibitem{arias}
J.~Arias de Reyna, \emph{An infinitely differentiable function with compact
support: Definition and properties}, English translation of the 1982 paper,
arXiv:1702.05442, 2017.
\url{https://arxiv.org/abs/1702.05442}.

\bibitem{renyi}
T.~van Erven and P.~Harremo\"es, \emph{R\'enyi divergence and
Kullback--Leibler divergence}, IEEE Transactions on Information Theory
\textbf{60} (2014), 3797--3820; arXiv:1206.2459.
\url{https://arxiv.org/abs/1206.2459}.

\bibitem{dlmf-beta}
NIST, \emph{Digital Library of Mathematical Functions}, \S5.12,
Euler's beta integral, especially Eq.~5.12.1. Accessed 30 September 2026.
\url{https://dlmf.nist.gov/5.12}.

\bibitem{dlmf-gamma}
NIST, \emph{Digital Library of Mathematical Functions}, \S5.11,
gamma and digamma asymptotics and gamma ratios. Accessed 30 September 2026.
\url{https://dlmf.nist.gov/5.11}.
\end{thebibliography}
\end{document}
